\chapter{Sicurezza su Internet}

\section{Secure E-mail e S/MIME}

S/MIME (Secure/Multipurpose Internet Mail Extension) è un'estensione di sicurezza allo standard di formattazione e-mail MIME. È stato progettato per fornire servizi crittografici avanzati ai messaggi di posta elettronica, tra cui autenticazione, integrità e riservatezza. 

\begin{notebox}[Base Crittografica di S/MIME]
S/MIME si basa fortemente sull’uso della crittografia a chiave pubblica asimmetrica e sull'infrastruttura a chiave pubblica (PKI), impiegando certificati digitali X.509 per la gestione sicura delle chiavi e per l'identificazione certa dei mittenti.
\end{notebox}

\subsection{S/MIME e Funzionalità}
La funzionalità S/MIME è oggi integrata in modo nativo nella maggior parte dei software di posta elettronica e aggiunge nuovi tipi di contenuto MIME per fornire la capacità di firmare e/o cifrare i messaggi, supportando quattro nuove funzioni specifiche:

\begin{description}
    \item[Dati Busta (Enveloped data):] Consiste in un contenuto cifrato di qualsiasi tipo e nelle chiavi simmetriche di cifratura del contenuto, a loro volta cifrate per uno o più destinatari. In questo modo solo i destinatari previsti, che sono in possesso della chiave privata corretta, possono decifrare la chiave e accedere al contenuto del messaggio.
    \item[Dati Firmati (Signed data):] Una firma digitale creata calcolando il message digest (hash) del contenuto e cifrandolo con la chiave privata del mittente. Il contenuto e la firma vengono codificati in base64. Un messaggio di questo tipo può essere visualizzato solo da un destinatario con capacità S/MIME e garantisce autenticità e integrità.
    \item[Dati Firmati in Chiaro (Clear-signed data):] Simile ai dati firmati, ma solo la firma digitale viene codificata in base64, mentre il testo rimane in chiaro. Permette ai destinatari privi di client con capacità S/MIME di leggere regolarmente il contenuto del messaggio (pur non potendo verificarne la firma), favorendo la compatibilità.
    \item[Dati Firmati e Busta (Signed and enveloped data):] Entità annidate in cui dati cifrati possono essere firmati, e dati firmati possono essere cifrati. Questa potente combinazione consente di ottenere simultaneamente riservatezza, autenticità e integrità in un unico messaggio strutturato.
\end{description}

\subsection{Funzionamento Generale}

\textbf{Fase di Firma (Dati Firmati/Firmati in Chiaro):}
\begin{enumerate}
    \item Viene calcolato un \textit{message hash} del contenuto (tipicamente utilizzando algoritmi robusti come SHA-256).
    \item L'hash viene cifrato con la chiave privata del mittente (es. tramite RSA o DSA) per creare la firma digitale vera e propria.
    \item La firma viene allegata al messaggio e codificata in base64 (una codifica che mappa i dati binari in caratteri ASCII stampabili sicuri) per proteggerla da alterazioni accidentali in transito.
\end{enumerate}

\textbf{Fase di Cifratura (Dati Busta):}
\begin{enumerate}
    \item S/MIME genera una chiave di sessione segreta pseudocasuale esclusivamente per il singolo messaggio.
    \item Il messaggio viene cifrato con questa chiave di sessione usando un algoritmo simmetrico rapido (es. AES).
    \item La chiave di sessione viene quindi cifrata con la chiave pubblica del destinatario (usando l'algoritmo RSA) e trasmessa come "busta" insieme al messaggio cifrato.
    \item Il destinatario userà la propria chiave privata per decifrare prima la chiave di sessione e, di conseguenza, il messaggio originale.
\end{enumerate}

\section{DomainKeys Identified Mail (DKIM)}

DKIM è una specifica per la firma crittografica dei messaggi e-mail che permette a un intero dominio mittente (es. \textit{@azienda.com}) di assumersi la responsabilità formale di un messaggio. I destinatari verificano la firma interrogando i record DNS pubblici del dominio mittente.

\subsection{Architettura della Posta su Internet}
Per comprendere DKIM è necessario definire i componenti:
\begin{itemize}
    \item \textbf{MUA (Message User Agent):} Il client di posta dell'utente finale (es. Outlook, Thunderbird).
    \item \textbf{MHS (Message Handling Service):} Composto da:
    \begin{itemize}
        \item \textit{MSA (Mail Submission Agent):} Accetta il messaggio iniziale dal MUA.
        \item \textit{MTA (Message Transfer Agent):} Il server che inoltra il messaggio nella rete.
        \item \textit{MDA (Mail Delivery Agent):} Il servizio finale che consegna il messaggio al Message Store.
    \end{itemize}
    \item \textbf{MS (Message Store):} Archivio messaggi in attesa di recupero (via protocolli POP/IMAP).
    \item \textbf{ADMD (Administrative Management Domain):} Provider di posta (ISP o organizzazione aziendale).
    \item \textbf{DNS:} Servizio usato da DKIM per pubblicare le chiavi pubbliche del dominio.
\end{itemize}

\subsection{Strategia di DKIM}

\begin{warnbox}[Differenza tra S/MIME e DKIM]
A differenza di S/MIME, che è incentrato sull'utente finale e protegge l'intero contenuto, DKIM è uno strumento implementato a \textbf{livello di server} (trasparente per l'utente) ed è fondamentale per combattere minacce come \textit{phishing} e \textit{spoofing} dell'indirizzo di posta.
\end{warnbox}

DKIM opera nel seguente modo:
\begin{enumerate}
    \item Il messaggio in uscita viene firmato dall'MTA con la chiave privata segreta del dominio amministrativo (ADMD), coprendo parti del contenuto e le intestazioni chiave. Non viene mai usata la chiave personale dell'utente.
    \item L'MDA o l'MTA del destinatario recupera la chiave pubblica del dominio mittente eseguendo una query al DNS (record TXT).
    \item Se la verifica ha successo, il server conferma che il messaggio proviene effettivamente dal dominio dichiarato e non è stato alterato in transito, migliorando la \textit{reputazione spam} del messaggio.
\end{enumerate}

\section{Secure Sockets Layer e Transport Layer Security}

TLS, il successore standardizzato di SSL, fornisce un servizio di sicurezza \textit{end-to-end} affidabile basandosi sul protocollo di trasporto TCP.

\subsection{Architettura TLS}
TLS è composto da un'architettura a due livelli principali, progettati in modo modulare.

\begin{defbox}[Connessione vs Sessione in TLS]
\begin{itemize}
    \item \textbf{Connessione TLS:} Una relazione transitoria e temporanea peer-to-peer per il trasporto sicuro dei dati. Ogni connessione è legata a una specifica sessione.
    \item \textbf{Sessione TLS:} Un'associazione logica tra client e server creata dall'Handshake Protocol. Definisce un insieme di parametri crittografici che possono essere condivisi tra molteplici connessioni simultanee o successive, evitando così le costose operazioni asimmetriche di rinegoziazione.
\end{itemize}
\end{defbox}

\subsection{Protocolli TLS in Dettaglio}
\begin{description}
    \item[Record Protocol:] Prende i dati dalle applicazioni, li frammenta, li comprime (opzionale), calcola un MAC (per garantirne l'integrità), li cifra (per la confidenzialità) e aggiunge un'intestazione per l'invio.
    \item[Change Cipher Spec Protocol:] Un singolo e breve messaggio per segnalare alla controparte il passaggio definitivo all'uso della nuova suite di cifratura appena negoziata.
    \item[Alert Protocol:] Invia avvisi strutturati (che possono essere errori fatali che chiudono la connessione, o semplici notifiche).
    \item[Handshake Protocol:] Il cuore di TLS. Autentica le parti, negozia gli algoritmi (cifratura simmetrica, hashing MAC e metodo di scambio chiavi) e stabilisce le chiavi segrete condivise. Si divide in 4 fasi principali (scambio Hello, invio certificati/chiavi dal server, invio certificati/chiavi dal client, e verifica finale con messaggi Finished).
    \item[Heartbeat Protocol:] Aggiunta recente per monitorare la disponibilità della controparte e prevenire la chiusura silente della connessione da parte dei firewall o dei NAT.
\end{description}
\section{HTTPS}

HTTPS (HTTP over SSL/TLS) combina il protocollo web HTTP e TLS per garantire comunicazioni sicure tra browser e server web (generalmente sulla porta 443). 

HTTPS si occupa di cifrare l'intera comunicazione applicativa, tra cui:
\begin{itemize}
    \item L'URL esatto del documento o della risorsa richiesta
    \item Il contenuto del documento trasferito e i dati dei moduli compilati dall'utente
    \item I cookie di sessione e tutte le intestazioni HTTP.
\end{itemize}

\section{Sicurezza di IPv4 e IPv6 (IPsec)}

IPsec (Internet Protocol Security) è una potente suite di protocolli che fornisce sicurezza direttamente a livello di rete (livello IP), garantendo autenticazione e cifratura del tutto trasparenti per le applicazioni o gli utenti. È lo standard per reti VPN aziendali, accessi remoti ed extranet. Se implementato direttamente su firewall o router di confine, protegge l'intero traffico perimetrale.

\subsection{Associazioni di Sicurezza (SA)}
Una \textbf{SA (Security Association)} è una relazione logica \textbf{unidirezionale} (simplex) tra mittente e destinatario che definisce esattamente quali servizi crittografici di sicurezza applicare al traffico che vi scorre. Per una comunicazione bidirezionale servono quindi due SA. Una SA è identificata univocamente da tre parametri: Indice dei Parametri di Sicurezza (SPI), l'indirizzo IP di destinazione e l'identificatore del protocollo di sicurezza (tipicamente ESP).

\subsection{Encapsulating Security Payload (ESP)}
ESP è il protocollo cardine di IPsec che fornisce confidenzialità (cifratura), integrità e autenticazione dell'origine dei dati (avendo di fatto sostituito l'obsoleto protocollo AH - Authentication Header). 

ESP supporta due distinte modalità operative:
\begin{description}
    \item[Modalità Transport:] Protegge solo il \textit{payload} del pacchetto IP (es. il singolo segmento TCP/UDP), ma lascia l'intestazione IP originale in chiaro. È usata prevalentemente per comunicazioni sicure end-to-end tra due singoli host.
    \item[Modalità Tunnel:] Protegge l'\textbf{intero pacchetto IP} originale, cifrandolo e incapsulandolo all'interno di un \textbf{nuovo pacchetto IP} dotato di una nuova intestazione esterna. È usata tipicamente per le VPN (da gateway a gateway), nascondendo completamente l'architettura della rete interna e il traffico ai router pubblici intermedi.
\end{description}
